
\section{Conclusions}
\label{sec:conclusions}
\engExpl{Describe the conclusions (reflect on the whole introduction given in Chapter 1).}


  
\engExpl{Discuss the positive effects and the drawbacks.\\
Describe the evaluation of the results of the degree project.\\
Did you meet your goals?\\
What insights have you gained?\\
What suggestions can you give to others working in this area?\\
If you had it to do again, what would you have done differently?}

\sweExpl{Uppfyllde du dina mål?\\
Vilka insikter har du fått?\\
Vilka förslag kan du ge till andra som arbetar inom detta område?
Om du skulle göra detta igen, vad skulle du ha gjort annorlunda?}

\section{Limitations}
\label{sec:limitations}
\engExpl{What did you find that limited your efforts? What are the limitations of your results?}

- Exploratory -> Fundamentally unclear whether a safe and lively version hybrid protocol with log syncing is even achievable

- Struggle with cloud reproducability hindering progress -> final analysis not as exhaustive as could have been

- Only considered if fast-path successful not if faster -> likely the main methological limitation on practicability study ->  guarantees on  end-to-end latency could be more impactful since less prior network layout study is necessary

- No stable storage -> Assumption was that this is not the dominating factor for end-to-end latency in global scenarios but there could be disk I/O overhead for the more complex adaptive version which needs to be considered

- Fixed learning rate of 0.005 for adaptive experiments -> basically just first value stumbled upon in literature \cite{Angelopoulus-NeuralInformationProcessing-2023} -> no study was made regarding its impact (unclear whether there is one) -> batch calibration used as a way to push threshold in the right range from where a low learning rate is acceptable

\section{Future work}
\label{sec:futureWork}

\engExpl{Describe valid future work that you or someone else could or should do.\\
Consider: What you have left undone? What are the next obvious things to be done? What hints can you give to the next person who is going to follow up on your work?}

\begin{itemize}
    \item Fully specify the hybrid protocol with safety and liveness proves to assess whether the approach actually works
    \item Find a way to give guarantees on actual end-to-end latency instead of success since a success does not always mean faster execution -> network latency changes might not impact success when contention is low, but make fast-path execution slower in comparison
    \item Determine a heuristic based on math instead of empirical data seems possible and logical -> worth exploring
    \item Try to integrate the approach into Nezha \cite{Geng-NezhaConsensus-2022}, which makes fast-path more viable
    \item => this work can give guidance on how to approach the problem and how to analyse it
\end{itemize}